\section{A uniform finite universe, including all empty cases}\label{sec:finite}

The preceding structural theory applies to infinite presentations. A second
useful result is more algorithmic: the entire finite theory can be made
uniform in numerical table codes. Here every verification quantifier is
bounded, and every loop invariant has a PR Boolean characteristic function.
This section supplies the constructions and their correctness proofs, rather
than assuming that an unspecified universal object can be found.

\subsection{The category of all coded finite sets}
Write $[n]=\{0,\ldots,n-1\}$, allowing $n=0$. Define
\[
 b=m+1,\qquad d_b(c,i)=\bigl\lfloor c/b^i\bigr\rfloor\bmod b.
\]
The divisor is positive, including when $m=0$. A table code for $[n]\to[m]$
is a number $c$ satisfying
\[
 c<(m+1)^n,\qquad\forall i<n\; d_{m+1}(c,i)<m.                 \tag{FT}
\]
The bounded formula (FT) is a PR predicate. Its interpretation is
$i\mapsto d_{m+1}(c,i)$. Conversely a bounded list $v_i<m$ is encoded by
$\sum_{i<n}v_i(m+1)^i$. The digit and reconstruction equations follow by
quantifier-free induction on the coded length, using division with remainder.
For $m=0$, a valid table exists exactly when $n=0$, and then its only code
is $0$. This convention eliminates a common hidden positivity assumption.

Identity is the code of $i\mapsto i$. For $f:[n]\to[m]$ and $g:[m]\to[k]$,
composition is the code of
\[
 i\longmapsto d_{k+1}\bigl(g,d_{m+1}(f,i)\bigr).
\]
The displayed bound on the inner digit gives the typing needed for the
outer table. Associativity follows at each input by associativity of
substitution of three table lookups, and table reconstruction then gives
literal equality of the codes. The same pointwise proof gives both unit
laws. The object predicate is true of every natural number; equality of
arrows is numerical equality. Hence this is a uniform iPRA presentation
$\mathsf{Fin}$, not a separate category chosen for each size bound.

The elementary rank/select tool will recur. Given a bit test $K(i)$ and
bound $B$, let
\[
 r(t)=\sum_{i<t}K(i),\quad \ell=r(B),\quad
 s(j)=\min\{i<B:K(i)=1\text{ and }r(i)=j\},                  \tag{RS}
\]
where the minimum returns $B$ if no witness exists. Then
\[
 \begin{split}
 j<\ell&\Rightarrow s(j)<B,\ K(s(j))=1,\ r(s(j))=j,\\
 i<B\wedge K(i)=1&\Rightarrow r(i)<\ell,\ s(r(i))=i.
 \end{split}                                               \tag{RS'}
\]
To prove these facts, scan $i=0,\ldots,B-1$: the counter increases by one
exactly at a successful entry. Its successive successful values are
$0,\ldots,\ell-1$, each attained once. The assertion for any current
prefix is a conjunction of bounded bit tests, so this is a legitimate
quantifier-free induction. There is no unbounded minimization.

\subsection{Initial objects, products, equalizers and all finite limits}
The object $[0]$ is initial: a table with empty domain has its unique code
$0$. The object $[1]$ is terminal: every entry must be $0$, so again there
is one table. The product of $[n]$ and $[m]$ is $[nm]$, with coordinates
\[
 \pi_1(t)=\lfloor t/m\rfloor,\qquad \pi_2(t)=t\bmod m.
\]
Division by zero is totalized by zero. If $nm=0$, the domain is empty, so
these defaults never enter a typing assertion. If $t<nm$, then $m>0$,
and the quotient and remainder equations give the required ranges.
Given $f:[z]\to[n]$, $g:[z]\to[m]$, the mediator is
$u(x)=m f(x)+g(x)$. Its projections are $f,g$, and any other mediator
must have the same quotient and remainder, hence the same entries.

For $f,g:[n]\to[m]$, filter $i<n$ by $f(i)=g(i)$. If $\ell$ entries
remain, the equalizer is $s:[\ell]\to[n]$ from (RS). A table
$h:[z]\to[n]$ with $fh=gh$ factors by $u(x)=r(h(x))$.
Equations (RS') prove $su=h$; applying $r$ to this equation proves
uniqueness. When no entry remains, $\ell=0$. A valid cone then has empty
domain, as follows from the same factorization proof.

\begin{theorem}[Uniform finite limits in $\mathsf{Fin}$]\label{thm:finite-limit}
There is a PR procedure which, from a finite table category $J$ and a
table functor $D:J\to\mathsf{Fin}$, constructs a limit, its projections,
and the unique mediating table of any supplied cone. The construction
and its universal property are provable in iPRA. Empty diagrams and
zero-valued objects are included.
\end{theorem}
\begin{proof}
Let the objects of $J$ be $0,\ldots,r-1$, and let $D(i)=[n_i]$. Enumerate
all tuples $x=(x_0,\ldots,x_{r-1})$ with $x_i<n_i$. A convenient
implementation enumerates words in radices $n_i+1$ and filters out the
spare digit in each coordinate. For $r=0$ there is one empty tuple; if
some $n_i=0$ and $r>0$, there are no tuples.
Retain exactly the tuples for which, for every arrow $a:i\to j$ of $J$,
\[
                         D(a)(x_i)=x_j.                     \tag{LC}
\]
This is a finite product of equality bits over the supplied arrow table,
not a parameter whose correctness is assumed. Let $L$ be the increasing
list of retained tuple codes and $\ell$ its length. Projection $p_i$ takes
an index in $L$ to coordinate $i$. Condition (LC) proves the cone laws.

Given a cone $q_i:[z]\to[n_i]$, at each $t<z$ form the tuple
$(q_i(t))_{i<r}$. It satisfies (LC), so (RS') assigns it a unique rank
$u(t)<\ell$. These ranks form the mediating table. Its projections are
$q_i$ by construction. If $v$ also has these projections, the tuple at
$v(t)$ equals the tuple at $u(t)$ coordinatewise; reconstruction and
injectivity of selection yield $v(t)=u(t)$.

Every loop ranges over a displayed finite code bound: objects, arrows,
tuple codes or cone inputs. The invariants say that a retained tuple
satisfies each earlier equation, or that the constructed output entries
have the required ranks and coordinates. All are bounded arithmetic
formulas and are Booleanized as in Lemma~\ref{lem:bounded}. This yields
uniform numerical iPRA proofs, including the input sizes as variables.
For the empty diagram, the sole empty tuple is retained and the limit
is $[1]$. If a diagram object is empty, its full product is empty and
the limit is $[0]$; a cone into it must itself have empty domain.
\end{proof}

\subsection{Finite quotients and all finite colimits}
For two tables $f,g:[p]\to[q]$, start with a label table
$\lambda_0(v)=v$ for $v<q$. At stage $i<p$ let
$a=\lambda_i(f(i))$ and $b=\lambda_i(g(i))$. Replace every label equal
to $b$ by $a$, leaving all other labels unchanged. If $a=b$ this does
nothing. Every label remains below $q$, and the label table remains
idempotent: its labels are precisely fixed points of the table.

The following stronger invariants explain why this merging does not
identify anything unjustified. After stage $i$:
\begin{enumerate}
\item the endpoints of every processed pair have equal labels;
\item if $h:[q]\to[z]$ equalizes all processed pairs, then
      $h(v)=h(\lambda_i(v))$ for every $v<q$;
\item each $v$ is connected to $\lambda_i(v)$ by a chain of processed
      pair identifications, used in either direction.
\end{enumerate}
For (2), suppose the next pair is equalized. The old invariant gives
$h(a)=h(f(i))=h(g(i))=h(b)$; replacing label $b$ by $a$ therefore
preserves $h(v)=h(\lambda(v))$. For (3), concatenate the two old paths
from $f(i),g(i)$ to $a,b$ with the new pair, and attach this path when
a label changes. If the old paths have length at most $B_i$, the path from $b$ to $a$
has length at most $2B_i+1$, and an updated path has length at most
$3B_i+1$. Thus $B_0=0$, $B_{i+1}=3B_i+1$ is an explicit PR bound.
A coded path table can be updated by the same primitive recursion.
For factorization only (2) is needed. Invariant (3), together with
(1) applied along a coded path, additionally proves that equality of
final labels is exactly the finite-path relation generated by the
pairs. This is a claim about a coded PR graph and its paths, not a
minimality assertion quantified over arbitrary arithmetic relations.

For invariant (2), fix the target size $z$ and candidate table $h$ as
numerical parameters during induction. The processed-pair condition and
the invariant both quantify only over bounded indices in coded tables.
Their Booleanized implication is therefore quantifier-free; only after
the induction do we generalize over $h,z$.

After all pairs, enumerate the fixed-point labels in increasing order,
say $e_0,\ldots,e_{k-1}$. Define
\[
 Q(v)=\text{the rank of }\lambda_p(v)\text{ among these labels}.
\]
For an equalizing table $h$ set $\bar h(j)=h(e_j)$. Invariant (2)
proves $\bar h Q=h$. Also $Q(e_j)=j$, so any factor has value $h(e_j)$
at $j$; this proves uniqueness. Thus $Q:[q]\to[k]$ is a coequalizer.
If $q=0$, validity of $f,g$ forces $p=0$; the empty label list gives
$k=0$, with exactly the same universal property. All invariants used
to obtain the finite-set universal property concern bounded tables
and their entries, so the induction remains quantifier-free after
Booleanization. It is not induction over arbitrary target categories.

\begin{theorem}[Uniform finite colimits]\label{thm:finite-colimit}
Every finite table diagram in $\mathsf{Fin}$ has a colimit constructed
by a PR algorithm, with a PR operation on cocones returning the unique
mediator. Its iPRA correctness includes empty diagrams and empty sets.
\end{theorem}
\begin{proof}
For $D:J\to\mathsf{Fin}$ form the tagged disjoint union
$U=\coprod_{i\in J}[n_i]$, a finite ordinal obtained by cumulative sums.
For each arrow $a:i\to j$ and $x<n_i$, create the pair
$(i,x),(j,D(a)x)$. The total number of pairs is
$\sum_{a\in\operatorname{Arr}J}n_{\operatorname{src}a}$.
Apply the preceding merging algorithm. Its quotient maps composed with
the tag injections form a cocone. A cocone $q_i:[n_i]\to[z]$
defines $h(i,x)=q_i(x)$ on $U$ and equalizes each generating pair.
The coequalizer factorization therefore gives its unique mediator.
For the empty diagram, $U$ is empty and the colimit is $[0]$.
\end{proof}

\subsection{Exponentials with \texorpdfstring{$0^0=1$}{0 to power 0 equals 1}}
Let $V_{n,m}$ be the increasing list of all codes satisfying (FT), and
write $q=|V_{n,m}|$. There is one choice for an empty list of entries;
appending one entry gives exactly $m$ choices. Therefore the counting
recurrence is $q(0,m)=1$, $q(n+1,m)=m q(n,m)$, whence
$q(n,m)=m^n$ by quantifier-free induction. In particular $0^0=1$,
while $0^n=0$ for $n>0$.

\begin{theorem}[Cartesian closure of coded finite sets]\label{thm:finite-ccc}
In $\mathsf{Fin}$ the object $[q]$, $q=m^n$, is an exponential
$[m]^{[n]}$. Evaluation and currying are uniformly PR, and their
universal equations have iPRA proofs for all $n,m,z$, including zero.
\end{theorem}
\begin{proof}
Let $s$ select a raw code from $V_{n,m}$ and $r$ rank a valid code.
Define $\operatorname{ev}(j,x)=d_{m+1}(s(j),x)$ for $j<q,x<n$.
This is a table on $[q]\times[n]$; a pair never exists if one factor
is empty. For $h:[z]\times[n]\to[m]$, form the row code
\[
 c_t=\sum_{x<n}h(t,x)(m+1)^x,
 \qquad\operatorname{curry}(h)(t)=r(c_t).
\]
Every row is valid, so the rank is below $q$. Digit reconstruction gives
$\operatorname{ev}(\operatorname{curry}(h)(t),x)=h(t,x)$.
Conversely, if $u:[z]\to[q]$, the row obtained by evaluating $u(t)$
is the raw table $s(u(t))$, whose rank is $u(t)$. Thus
\[
 \operatorname{ev}(\operatorname{curry}h\times\id)=h,
 \qquad\operatorname{curry}(\operatorname{ev}(u\times\id))=u.
\]
The first is the factorization law and the second proves uniqueness.
If $n=0$, every row is empty and $q=1$. If $m=0<n$, then $q=0$;
a table $h:[z]\times[n]\to[0]$ forces $z=0$, so its transpose
has the required empty domain. No positivity premise has been omitted.
\end{proof}

\subsection{Finite functor categories and genuine finite Yoneda bijections}
A finite table category consists of finite object and arrow tables,
source, target, identity and composition entries, with a bounded checker
for all typing, unit and associativity equations. Given such $C,D$,
enumerate all assignments on their finite object and arrow domains and
retain precisely those satisfying the functor laws. Likewise enumerate
all object-indexed arrow tuples and retain the natural transformations.
Pointwise identities and vertical composition make a finite table category
$[C,D]$. The verification is the same finite naturality calculation as
in Section~\ref{sec:yoneda}, now with bounded quantifiers on actual tables.

For a finite $\mathsf{Fin}$-valued presheaf $P$, whose fibres are
finite ordinals with numerical equality, all transformations
$h_A\Rightarrow P$ can be enumerated in this manner. The Yoneda
constructions become reciprocal numerical maps between two finite lists.
Consequently $|\operatorname{Nat}(h_A,P)|=|P(A)|$, with an explicit
bijection, rather than merely a schematic correspondence.

Finite adjunction existence is also decidable by bounded enumeration:
enumerate candidate right-adjoint functor tables, then unit and counit
tables, and check naturality and the two triangles. This does not assert
that a right adjoint always exists. It returns either displayed adjunction
data or a finite certificate that no candidate table satisfies the laws.
The bounds may be enormous; primitive recursive arithmetic imposes no
polynomial-time requirement.
